% GENERATED FILE. Edit canonical lessons, then run npm run generate:books.
% canonical-source-hash: fnv1a64:977d93ea566e786b
% canonical-lessons: PA-C288-dhupp, PA-C288-lu, PA-C288-tapna, PA-C288-hunal, PA-C288-soka

\chapter{The Hot Season Ahead}
\label{ch:pa-a2-the-hot-season-ahead}
\hlchaptermodality{288}

\begin{chapteropening}
\textbf{By the end of this chapter:} \emph{I can say what the hot months will be like.}

Complete the last lesson of chapter 288: I can say what the hot months will be like.
\end{chapteropening}

\section[\texorpdfstring{\pa{ਧੁੱਪ} (dhupp)}{ਧੁੱਪ (dhupp)}]{\pa{ਧੁੱਪ} (dhupp) — sunshine}

\label{lesson:PA-C288-dhupp}



\begin{quote}
Before the new one: say the Punjabi for to drip, then the Punjabi for mud.
\end{quote}

\subsection*{You'll want to know: \pa{ਧੁੱਪ}}
\textbf{\pa{ਧੁੱਪ}} — \emph{dhupp} — \textquotedblleft{}sunshine\textquotedblright{}.

\begin{quote}
\pa{ਕੱਲ੍ਹ} \pa{ਤੇਜ਼} \pa{ਧੁੱਪ} \pa{ਹੋਵੇਗੀ।}
\end{quote}

\begin{quote}
\emph{kall tez dhupp hovegī.}
\end{quote}

\begin{quote}
Tomorrow there will be strong sunshine.
\end{quote}

A feminine noun.

\begin{cousinweb}
\textbf{\pa{ਧੁੱਪ}} (\emph{dhupp}) is inherited from a reconstructed \emph{*dhuppā}, \textquotedblleft{}sunshine\textquotedblright{}; CDIAL lists Punjabi \emph{dhupp} under it (CDIAL 6825).
\end{cousinweb}

\begin{grammarlens}[title={What you've built}]
One more word for saying what comes next.
\end{grammarlens}

\subsection*{Your turn}
\begin{itemize}
\raggedright
  \item \emph{Say it:} \emph{dhupp}
  \item \emph{Say it:} \emph{dhupp}, once more
  \item \emph{Say it:} \emph{kall tez dhupp hovegī.} — Tomorrow there will be strong sunshine.
  \item \emph{Recall:} say the Punjabi for to drip, then the Punjabi for mud, then say \emph{dhupp} again
\end{itemize}

\begin{tcolorbox}[breakable,colback=teal!4,colframe=teal!35!black,title={Before you move on}]
What does \pa{ਧੁੱਪ} mean? (\textquotedblleft{}Sunshine\textquotedblright{}.) Say it once more.
\end{tcolorbox}
\section[\texorpdfstring{\pa{ਲੂ} (lū)}{ਲੂ (lū)}]{\pa{ਲੂ} (lū) — a hot wind; the summer heatwave}

\label{lesson:PA-C288-lu}



\begin{quote}
Before the new one: say the Punjabi for mud, then the Punjabi for sunshine.
\end{quote}

\subsection*{You'll want to know: \pa{ਲੂ}}
\textbf{\pa{ਲੂ}} — \emph{lū} — \textquotedblleft{}a hot wind; the summer heatwave\textquotedblright{}.

\begin{quote}
\pa{ਲੂ} \pa{ਵਿੱਚ} \pa{ਅਸੀਂ} \pa{ਬਾਹਰ} \pa{ਨਹੀਂ} \pa{ਜਾਂਦੇ} \pa{ਹਾਂ।}
\end{quote}

\begin{quote}
\emph{lū vich asīṁ bāhar nahīṁ jānde hāṁ.}
\end{quote}

\begin{quote}
We do not go out in the hot wind.
\end{quote}

A feminine noun.

\begin{cousinweb}
\textbf{\pa{ਲੂ}} (\emph{lū}) goes back to a reconstructed \emph{*lūṣā}, \textquotedblleft{}burning\textquotedblright{}; CDIAL lists Punjabi only \emph{lūh}, \textquotedblleft{}burnt, scalded\textquotedblright{}, under it; the sense \textquotedblleft{}hot wind\textquotedblright{} is the one it gives for Lahnda \emph{lō} (CDIAL 11099).
\end{cousinweb}

\begin{grammarlens}[title={What you've built}]
One more word for saying what comes next.
\end{grammarlens}

\subsection*{Your turn}
\begin{itemize}
\raggedright
  \item \emph{Say it:} \emph{lū}
  \item \emph{Say it:} \emph{lū}, once more
  \item \emph{Say it:} \emph{lū vich asīṁ bāhar nahīṁ jānde hāṁ.} — We do not go out in the hot wind.
  \item \emph{Recall:} say the Punjabi for mud, then the Punjabi for sunshine, then say \emph{lū} again
  \item \emph{Write it:} copy \textbf{\pa{ਲੂ}} once, checking the whole word
\end{itemize}

\begin{tcolorbox}[breakable,colback=teal!4,colframe=teal!35!black,title={Before you move on}]
What does \pa{ਲੂ} mean? (\textquotedblleft{}A hot wind; the summer heatwave\textquotedblright{}.) Say it once more.
\end{tcolorbox}
\section[\texorpdfstring{\pa{ਤਪਣਾ} (tapṇā)}{ਤਪਣਾ (tapṇā)}]{\pa{ਤਪਣਾ} (tapṇā) — to blaze with heat; to become very hot}

\label{lesson:PA-C288-tapna}



\begin{quote}
Before the new one: say the Punjabi for sunshine, then the Punjabi for a hot wind.
\end{quote}

\subsection*{You'll want to know: \pa{ਤਪਣਾ}}
\textbf{\pa{ਤਪਣਾ}} — \emph{tapṇā} — \textquotedblleft{}to blaze with heat; to become very hot\textquotedblright{}.

\begin{quote}
\pa{ਜੂਨ} \pa{ਵਿੱਚ} \pa{ਧਰਤੀ} \pa{ਤਪੇਗੀ।}
\end{quote}

\begin{quote}
\emph{jūn vich dhartī tapegī.}
\end{quote}

\begin{quote}
In June the earth will blaze with heat.
\end{quote}

Past (it …): \pa{ਉਹ} \pa{ਪੱਥਰ} \pa{ਤਪਿਆ}, \emph{oh patthar tapiā}.

\begin{cousinweb}
\textbf{\pa{ਤਪਣਾ}} (\emph{tapṇā}) is inherited from Sanskrit \emph{tapyati}, \textquotedblleft{}is hot\textquotedblright{}; CDIAL lists Punjabi \emph{tapṇā}, \textquotedblleft{}to be heated\textquotedblright{}, under it (CDIAL 5684).
\end{cousinweb}

\begin{grammarlens}[title={What you've built}]
One more word for saying what comes next.
\end{grammarlens}

\subsection*{Your turn}
\begin{itemize}
\raggedright
  \item \emph{Say it:} \emph{tapṇā}
  \item \emph{Say it:} \emph{tapṇā}, once more
  \item \emph{Say it:} \emph{jūn vich dhartī tapegī.} — In June the earth will blaze with heat.
  \item \emph{Recall:} say the Punjabi for sunshine, then the Punjabi for a hot wind, then say \emph{tapṇā} again
  \item \emph{Write it:} copy \textbf{\pa{ਤਪਣਾ}} once, checking the whole word
\end{itemize}

\begin{tcolorbox}[breakable,colback=teal!4,colframe=teal!35!black,title={Before you move on}]
What does \pa{ਤਪਣਾ} mean? (\textquotedblleft{}To blaze with heat; to become very hot\textquotedblright{}.) Say it once more.
\end{tcolorbox}
\section[\texorpdfstring{\pa{ਹੁਨਾਲ} (hunāl)}{ਹੁਨਾਲ (hunāl)}]{\pa{ਹੁਨਾਲ} (hunāl) — the hot season}

\label{lesson:PA-C288-hunal}



\begin{quote}
Before the new one: say the Punjabi for a hot wind, then the Punjabi for to blaze with heat.
\end{quote}

\subsection*{You'll want to know: \pa{ਹੁਨਾਲ}}
\textbf{\pa{ਹੁਨਾਲ}} — \emph{hunāl} — \textquotedblleft{}the hot season\textquotedblright{}.

\begin{quote}
\pa{ਹੁਨਾਲ} \pa{ਵਿੱਚ} \pa{ਅਸੀਂ} \pa{ਛਾਂ} \pa{ਵਿੱਚ} \pa{ਬੈਠਦੇ} \pa{ਹਾਂ।}
\end{quote}

\begin{quote}
\emph{hunāl vich asīṁ chhāṁ vich baiṭhde hāṁ.}
\end{quote}

\begin{quote}
In the hot season we sit in the shade.
\end{quote}

A masculine noun.

\begin{cousinweb}
\textbf{\pa{ਹੁਨਾਲ}} (\emph{hunāl}) is inherited from Sanskrit \emph{uṣṇakāla}, \textquotedblleft{}hot season\textquotedblright{}; CDIAL lists Punjabi \emph{hunālā} under it (CDIAL 2391).
\end{cousinweb}

\begin{grammarlens}[title={What you've built}]
One more word for saying what comes next.
\end{grammarlens}

\subsection*{Your turn}
\begin{itemize}
\raggedright
  \item \emph{Say it:} \emph{hunāl}
  \item \emph{Say it:} \emph{hunāl}, once more
  \item \emph{Say it:} \emph{hunāl vich asīṁ chhāṁ vich baiṭhde hāṁ.} — In the hot season we sit in the shade.
  \item \emph{Recall:} say the Punjabi for a hot wind, then the Punjabi for to blaze with heat, then say \emph{hunāl} again
  \item \emph{Write it:} copy \textbf{\pa{ਹੁਨਾਲ}} once, checking the whole word
\end{itemize}

\begin{tcolorbox}[breakable,colback=teal!4,colframe=teal!35!black,title={Before you move on}]
What does \pa{ਹੁਨਾਲ} mean? (\textquotedblleft{}The hot season\textquotedblright{}.) Say it once more.
\end{tcolorbox}
\section[\texorpdfstring{\pa{ਸੋਕਾ} (sokā)}{ਸੋਕਾ (sokā)}]{\pa{ਸੋਕਾ} (sokā) — a drought}

\label{lesson:PA-C288-soka}



\begin{quote}
Before the new one: say the Punjabi for to blaze with heat, then the Punjabi for the hot season.
\end{quote}

\subsection*{You'll want to know: \pa{ਸੋਕਾ}}
\textbf{\pa{ਸੋਕਾ}} — \emph{sokā} — \textquotedblleft{}a drought\textquotedblright{}.

\begin{quote}
\pa{ਜੇ} \pa{ਮੀਂਹ} \pa{ਨਹੀਂ} \pa{ਆਵੇਗਾ} \pa{ਤਾਂ} \pa{ਸੋਕਾ} \pa{ਹੋਵੇਗਾ।}
\end{quote}

\begin{quote}
\emph{je mīṁh nahīṁ āvegā tāṁ sokā hovegā.}
\end{quote}

\begin{quote}
If the rain does not come, there will be a drought.
\end{quote}

A masculine noun.

\begin{cousinweb}
\textbf{\pa{ਸੋਕਾ}} (\emph{sokā}) is inherited from a reconstructed \emph{*śauṣkya}, \textquotedblleft{}dryness\textquotedblright{}; CDIAL lists Punjabi \emph{sok} and \emph{sokkā}, \textquotedblleft{}dryness, drought\textquotedblright{}, under it (CDIAL 12652).
\end{cousinweb}

\begin{grammarlens}[title={What you've built}]
One more word for saying what comes next.
\end{grammarlens}

\subsection*{Your turn}
\begin{itemize}
\raggedright
  \item \emph{Say it:} \emph{sokā}
  \item \emph{Say it:} \emph{sokā}, once more
  \item \emph{Say it:} \emph{je mīṁh nahīṁ āvegā tāṁ sokā hovegā.} — If the rain does not come, there will be a drought.
  \item \emph{Recall:} say the Punjabi for to blaze with heat, then the Punjabi for the hot season, then say \emph{sokā} again
  \item \emph{Write it:} copy \textbf{\pa{ਸੋਕਾ}} once, checking the whole word
\end{itemize}

\begin{tcolorbox}[breakable,colback=teal!4,colframe=teal!35!black,title={Before you move on}]
What does \pa{ਸੋਕਾ} mean? (\textquotedblleft{}A drought\textquotedblright{}.) Say it once more.
\end{tcolorbox}
